\documentclass[11pt]{article}
\usepackage[utf8]{inputenc}
\usepackage[T1]{fontenc}
\usepackage{lmodern}
\usepackage[margin=0.9in]{geometry}
\usepackage{parskip}
\usepackage[hidelinks]{hyperref}
\pagestyle{empty}
\setlength{\parskip}{6pt}
\linespread{1.03}

\begin{document}

\begin{center}
{\large\bfseries New Kids on the Block? The Effects of Neighborhood Childcare Expansion\\[2pt] on Maternal Employment and Fertility}

\vspace{8pt}
Pedro Carneiro$^{1}$ \quad Andr\'es Hojman$^{2}$ \quad Micaella Magnasco$^{2}$ \quad Marigen Narea$^{2}$ \quad Jorge Rodr\'iguez$^{3}$

\vspace{3pt}
{\small $^{1}$University College London \quad $^{2}$Pontificia Universidad Cat\'olica de Chile \quad $^{3}$Universidad de los Andes, Chile}
\end{center}

\vspace{2pt}
\noindent\textbf{JEL codes:} J13, J22\\
\textbf{Keywords:} early childhood education; maternal employment; staggered difference-in-differences

\vspace{2pt}
\hrule
\vspace{4pt}

\textbf{Research question.} Female labor force participation in Chile has stalled at around 50\%, and the total fertility rate has fallen to 1.17, one of the lowest in the world (OECD, 2025). Between 2015 and 2024, Chile substantially expanded public childcare and preschool provision for children aged 0--4 through new public centers. We ask whether opening a public early childhood education center in a neighborhood increases mothers' formal employment and earnings, and whether it affects fertility. The question is directly relevant to the ongoing debate on universal childcare coverage and the design of National Care Systems.

\textbf{Prior work and contribution.} Access to childcare can raise maternal labor supply and, in some settings, fertility (Baker, Gruber and Milligan, 2008; Havnes and Mogstad, 2011; Hojman and Lopez Boo, 2022; Attanasio et al., 2024). Causal evidence for Latin America, however, remains scarce. Moreover, almost all studies of preschool expansions rely on municipal, regional or aggregate variation. To our knowledge, this is the first study to combine geocoded center openings with individual-level administrative data at the level of the neighborhood unit (\emph{unidad vecinal}). This geography is much finer than the municipality, which sharpens the measure of local access.

\textbf{Empirical strategy.} We exploit the staggered opening of the first public center across neighborhood units. In an event-study design, we compare units that receive their first center with units that never do. We use heterogeneity-robust estimators that avoid the known biases of two-way fixed effects regressions under staggered adoption (Goodman-Bacon, 2021; Baker et al., 2026). The estimates are intention-to-treat effects of local access. Treated and never-treated units are similar in the characteristics of the women who live there and differ mainly in population size, a level difference absorbed by unit fixed effects. Results are robust to a cleaner control group that excludes never-treated units exposed to openings in neighboring units.

\textbf{Data.} We assembled a geocoded panel of all public early childhood education centers from the providers' administrative records. We link it to a secure data environment where Chile's Ministry of Social Development gives competitively selected research teams access to linked administrative records. For women in the Social Household Registry, we observe preschool enrollment, Civil Registry births, and formal employment and earnings from tax, unemployment insurance and pension records.

\textbf{Preliminary findings.} Our results show a clear first stage: enrollment rises in the year a center opens and stays higher, with no pre-trends. Mothers' formal employment and earnings increase significantly in the years after the opening. The earnings gain comes from more mothers working, not from higher wages among those already employed. We find no effect on fertility.

\textbf{Limitations and further work.} Further work will adjust for non-compliance and study the mechanisms behind these effects, including which families use the new centers and whether they replace other care arrangements. It will also study heterogeneity across mothers and neighborhoods, and longer-run outcomes for mothers and children.

\end{document}
